\documentclass[11pt]{article}
\usepackage[a4paper,margin=26mm]{geometry}
\usepackage{amsmath,amssymb}
\setlength{\parindent}{0pt}
\setlength{\parskip}{7pt}
\emergencystretch=2em
\title{A star need not maximize graph spectral radius\\
\large Disproof of conjecture 00000007164}
\author{ziangni-sys}
\date{4 October 2026}
\begin{document}
\maketitle
The English conjecture says that the maximal-radius graph is the star.
The Chinese version likewise asserts star-type maximal radius.
There is no restriction to trees or to a fixed number of edges.
We compare two actual connected simple graphs on the same three vertices:
the star $K_{1,2}$ and the complete graph $K_3$.
The latter has strictly greater adjacency spectral radius.

\textbf{The graphs and actual adjacency matrices.}
Let the vertex set be $\{0,1,2\}$. The star has exactly the edges
$\{0,1\}$ and $\{0,2\}$. The triangle has those edges and also $\{1,2\}$.
They are both simple and connected: vertex $0$ is adjacent to each
other vertex, providing paths between every pair. Their adjacency matrices
over $\mathbb C$ are
\[
 S=\begin{pmatrix}0&1&1\\1&0&0\\1&0&0\end{pmatrix},
 \qquad
 T=\begin{pmatrix}0&1&1\\1&0&1\\1&1&0\end{pmatrix}.
\]
The star has two edges and the triangle three. Thus this counterexample
uses a fixed vertex count, not a fixed edge count. It makes no claim
against extremality of stars within the class of trees or another
restricted family.

\textbf{Full actual spectra.}
For a finite complex matrix $M$, its spectrum consists of those $z$
for which $zI-M$ is not invertible. This is equivalent to
$\det(zI-M)=0$, since a square matrix over a field is invertible exactly
when its determinant is nonzero.

Direct expansion of the actual adjacency determinants gives
\[
 \det(zI-S)=z(z^2-2),\qquad
 \det(zI-T)=(z-2)(z+1)^2.
\]
Writing $r=\sqrt2>0$, we have $r^2=2$ and
$z^2-2=(z-r)(z+r)$. Therefore these identities characterize every
complex spectral point:
\[
 \sigma(S)=\{0,\sqrt2,-\sqrt2\},\qquad
 \sigma(T)=\{2,-1\}.
\]
In particular these are complete spectra, not just a selected pair of
eigenvalues or assumed scalar certificates.

\newpage
\textbf{Spectral radii and strict comparison.}
The standard finite-dimensional spectral radius is
\[
 \rho(M)=\sup\{|z|:z\in\sigma(M)\}.
\]
Using the actual spectra computed above gives
\[
 \rho(S)=\sqrt2,\qquad \rho(T)=2.
\]
Since $\sqrt2\geq0$ and $(\sqrt2)^2=2<4=2^2$, we have
\[
 \rho(S)=\sqrt2<2=\rho(T).
\]
Consequently the star is not a maximal-radius graph even among
connected simple graphs on three fixed vertices.

This falsifies the first unqualified conjunct of the conjecture.
It does not require interpreting the separate statement about the
neighborhood of a star, nor does it impose or contradict a theorem
with an additional fixed-edge or tree hypothesis. Such restrictions
would change the extremal problem, and are absent from the source.

\textbf{Formal correspondence.}
The Lean source constructs the star and triangle as actual
\texttt{SimpleGraph (Fin 3)} objects. Their symmetry and absence of
loops are verified, and \texttt{star\_connected} and
\texttt{triangle\_connected} prove actual graph connectivity.
The finite set of ordered endpoint representatives $i<j$ of undirected
edges is defined from adjacency. \texttt{edge\_counts} proves its
cardinalities are two and three.

All matrix computations use the graphs' actual
\texttt{SimpleGraph.adjMatrix} values over $\mathbb C$.
\texttt{spectrum\_det} connects Mathlib's actual \texttt{spectrum}
to determinant vanishing via matrix invertibility.
\texttt{star\_determinant} and \texttt{triangle\_determinant} prove
the determinant identities for every complex parameter.
\texttt{star\_spectrum} and \texttt{triangle\_spectrum} then prove
equality of the full actual spectrum sets.

The definition \texttt{radius} is the real supremum of the moduli
of the actual complex spectrum.
\texttt{star\_radius} and \texttt{triangle\_radius} evaluate these
suprema; \texttt{radius\_strict} proves their strict comparison.
\texttt{star\_not\_maximal} directly negates the maximal-radius
property quantified over simple graphs on the same three vertices.
The final \texttt{counterexample} combines the graph hypotheses,
actual spectra and strict failure of maximality.

\textbf{Reproduction.}
Use Lean 4.19.0 with public Mathlib revision
\begin{center}\small\texttt{c44e0c8ee63ca166450922a373c7409c5d26b00b}\end{center}
Run \texttt{lake build} in \texttt{lean/}. The project prints axiom audits.
The validation notes record the final whole-project build and PDF inspection.
\end{document}
